% Documentation of the persian-calendar package. Build: make doc   (XeLaTeX, font Parisa)
\documentclass[11pt]{article}

\usepackage[a4paper,margin=2.5cm]{geometry}
\usepackage{persian-calendar}
\usepackage{pgffor}
\usepackage[colorlinks,linkcolor=blue!50!black,urlcolor=blue!50!black]{hyperref}
\usepackage{xepersian}

\settextfont{Parisa}
\setlatintextfont[Extension=.otf,BoldFont=lmroman10-bold,ItalicFont=lmroman10-italic,BoldItalicFont=lmroman10-bolditalic]{lmroman10-regular}

% --- helpers --------------------------------------------------------------
% \macro|syntax|: heading line for a macro description
\NewDocumentCommand\macro{v}{\par\bigskip\noindent\hbox{\texttt{\bfseries #1}}\par\nopagebreak\smallskip}
% \ex|code|{result}: a call, an arrow and its typeset result (the call is typed twice on purpose)
\NewDocumentCommand\ex{v m}{\par\noindent\hspace*{1em}\texttt{#1}\quad$\Rightarrow$\quad#2}
\newcommand\iso[1]{\lr{#1}}            % numeric dates must stay left to right
\newcommand\pkg[1]{\textsf{#1}}
\newcommand\cs[1]{\texttt{\textbackslash #1}}
\newcommand\gap{\par\smallskip}
\newcount\jd \newcount\wd
\renewcommand\thepage{\lr{\arabic{page}}}   % page numbers in Latin digits (the footer is typeset in the Persian font)
% table cells for the name lists: \mc{n} month n, \yc{y} year y; \MROW/\YROW{first} fill a row of four
\newcommand\mc[1]{\edef\mm{\number\numexpr#1\relax}\texttt{\mm}\quad\pcalendarmonthname{\mm}}
\newcommand\yc[1]{\edef\yy{\number\numexpr#1\relax}\texttt{\yy}\quad\pcalendaryearname{\yy}}
\newcommand\MROW[1]{\mc{#1}&\mc{#1+3}&\mc{#1+6}&\mc{#1+9}}
\newcommand\YROW[1]{\yc{#1}&\yc{#1+3}&\yc{#1+6}&\yc{#1+9}}

% The prose uses the Latin font, so the macros that print Persian text switch to the Persian font.
\def\fa@wrap#1{\expandafter\let\csname fa@orig@\string#1\endcsname#1%
  \def#1##1{{\setpersianfont\csname fa@orig@\string#1\endcsname{##1}}}}
\fa@wrap\pdate \fa@wrap\gdate \fa@wrap\pcalendarmonthname \fa@wrap\pcalendargregorianmonthname
\fa@wrap\pcalendarweekdayname \fa@wrap\pcalendarweekdayshortname \fa@wrap\pcalendaryearname

\title{The \pkg{persian-calendar} package}
\author{Javad Razavian\\\texttt{javadr@gmail.com}}
\date{Version \pcalendarversion\quad(\pcalendardate)}

\begin{document}
\latin % the text is English (Latin font, left to right); the typeset results are Persian

\maketitle

\begin{abstract}\noindent
\pkg{persian-calendar} converts dates between the Persian (Jalali, Solar Hijri) and Gregorian calendars
inside \LaTeX{}. It prints dates with Persian month names, gives day numbers you can do arithmetic
on, and provides Persian month, weekday and year-animal names. It is built on the \pkg{pgfcalendar}
package. The conversions work with any engine; printing the Persian names needs a Unicode
engine and a Persian font, for example XeLaTeX with \pkg{xepersian}.
\end{abstract}

\tableofcontents

% ---------------------------------------------------------------------------
\section{What the package does and does not do}\label{sec:what1}

The package is a date converter and a set of name tables. It provides:
\begin{itemize}
  \item Gregorian $\to$ Jalali and Jalali $\to$ Gregorian conversion, printed as text or as numbers;
  \item Julian day numbers for Jalali dates, so you can add days, subtract dates or loop over days;
  \item Persian names of the Jalali and Gregorian months and of the weekdays;
  \item the animal name of a year in the twelve-animal cycle.
\end{itemize}
It does \emph{not} typeset a month grid or a wall calendar, it has no holidays, and it has no macro
that returns the weekday of a date (section~\ref{sec:names6} shows how to get one in three lines).

\paragraph{Requirements.} The \pkg{pgf} bundle (which contains \pkg{pgfcalendar}) is all the
conversions need, and they work with pdf\LaTeX{}, Xe\LaTeX{} and Lua\LaTeX{}. Only the names (and
so \cs{pdate}, \cs{gdate} and the \texttt{\textbackslash pcalendar\ldots name} macros) are Persian
text: they are stored as UTF-8 in the source file, so printing them needs an engine and packages that
handle UTF-8 Persian letters, together with a font that has them. \pkg{xepersian} with Xe\LaTeX{} is
what this manual uses, but any setup that can typeset Persian will do. Without one, use the numeric
macros \cs{pisodate} and \cs{gisodate}, which print only digits.

% ---------------------------------------------------------------------------
\section{Installation}\label{sec:installation2}

From a TeX distribution that includes the package, nothing needs doing. To install it by hand, copy
\texttt{persian-calendar.sty} into a directory that TeX searches, for example the local tree
\texttt{texmf/tex/latex/persian-calendar/}, and refresh the file database with \texttt{mktexlsr}. In the
source repository, \texttt{make dev} does this for you. You can also keep \texttt{persian-calendar.sty}
next to your document.

Load the package in the preamble. The examples here use \pkg{xepersian} for the Persian text and load
\pkg{persian-calendar} before it:
\begin{verbatim}
\usepackage{persian-calendar}
\usepackage{xepersian}
\settextfont{Parisa}
\end{verbatim}
The package takes no options.

% ---------------------------------------------------------------------------
\section{Quick start}\label{sec:quick3}

\begin{verbatim}
\pdate{2015-03-21}          % Gregorian date in, Jalali text out
\gdate{1394-01-01}          % Jalali date in, Gregorian text out
\pisodate{2015-03-21}       % Gregorian in, Jalali as DD/MM/YYYY
\gisodate{1394-01-01}       % Jalali in, Gregorian as DD/MM/YYYY
\end{verbatim}
which print:
\ex|\pdate{2015-03-21}|{\pdate{2015-03-21}}
\ex|\gdate{1394-01-01}|{\gdate{1394-01-01}}
\ex|\pisodate{2015-03-21}|{\iso{\pisodate{2015-03-21}}}
\ex|\gisodate{1394-01-01}|{\iso{\gisodate{1394-01-01}}}
\gap
Today's date, in the Jalali calendar, is \cs{pdate}\texttt{\{\cs{year}-\cs{month}-\cs{day}\}}, which
prints \pdate{\year-\month-\day}.

% ---------------------------------------------------------------------------
\section{Date conversion}\label{sec:date4}

All four macros take one argument, a date written \texttt{Y-M-D}: a year, a month number and a day
number joined by hyphens. The month and day may have a leading zero or not
(\texttt{2015-3-21} and \texttt{2015-03-21} are the same).

\macro|\pdate{Y-M-D}|
The argument is a \emph{Gregorian} date. Prints the \emph{Jalali} date as ``day month-name year'',
for example \pdate{2024-03-20}.

\macro|\gdate{Y-M-D}|
The argument is a \emph{Jalali} date. Prints the \emph{Gregorian} date as ``day month-name year'',
using the Persian transliteration of the Gregorian month, for example \gdate{1403-01-01}.

\macro|\pisodate{Y-M-D}|
Like \cs{pdate} (Gregorian in, Jalali out) but prints numbers: day, month and year, each day and
month padded to two digits and separated by a slash, so \iso{\pisodate{2024-03-20}}. Note the order:
this is day/month/year, not the ISO order year-month-day.

\macro|\gisodate{Y-M-D}|
Like \cs{gdate} (Jalali in, Gregorian out) but prints numbers in the same day/month/year form, for
example \iso{\gisodate{1403-01-01}}.

\paragraph{Leap years.} Esfand, the last Jalali month, has 30 days in a leap year:
\ex|\pdate{2025-03-20}|{\pdate{2025-03-20}}
\ex|\gdate{1403-12-30}|{\gdate{1403-12-30}}
\ex|\gdate{1404-12-30}|{\gdate{1404-12-30}}
The last line is an invalid date, since 1404 is not a leap year; the package does not check this
(section~\ref{sec:limitations9}).

\paragraph{Wrap the numeric forms.} Digits in Persian text are reordered by the bidirectional
algorithm. To keep \cs{pisodate} and \cs{gisodate} in day/month/year order inside Persian text, put
them in an \cs{lr} box, as this document does.

% ---------------------------------------------------------------------------
\section{Day numbers and date arithmetic}\label{sec:day5}

The conversion macros above cannot add days. For that, the package works with Julian day numbers as
defined by \pkg{pgfcalendar}: whole days counted from a fixed origin, so subtracting two numbers gives
the number of days between two dates.

\macro|\pcalendardatetojulian{Y-M-D}{\count}|
Stores the day number of the \emph{Jalali} date \texttt{Y-M-D} in the count register \cs{count}. Append
\texttt{+N} to the date to add \texttt{N} days first: \texttt{1394-01-01+10} is ten days after
1 Farvardin 1394.

\macro|\pcalendarjuliantodate{\count}{\y}{\m}{\d}|
Converts the day number in \cs{count} back to a \emph{Jalali} date and puts the year, month and day
into the macros \cs{y}, \cs{m} and \cs{d} (you may choose other names).

\begin{verbatim}
\newcount\jd
\pcalendardatetojulian{1394-01-01+10}\jd
\pcalendarjuliantodate\jd\y\m\d
\number\y-\number\m-\number\d
\end{verbatim}
\pcalendardatetojulian{1394-01-01+10}\jd%
\pcalendarjuliantodate\jd\y\m\d%
prints \iso{\number\y-\number\m-\number\d}.

\paragraph{Read \cs{y}, \cs{m}, \cs{d} immediately.} After \cs{pcalendarjuliantodate} the three macros
are aliases of the package's internal result registers, not copies. Any later conversion
(\cs{pdate}, \cs{gisodate}, and so on) changes them. If you need the values later, store them in your
own counters first: \texttt{\cs{newcount}\cs{myyear}\ \cs{myyear=\cs{y}}}.

\paragraph{Offsets across a year end.} Day arithmetic knows about leap years:
\pcalendardatetojulian{1403-12-29+1}\jd\pcalendarjuliantodate\jd\y\m\d%
1403-12-29 plus one day is \iso{\number\y-\number\m-\number\d}, because 1403 is a leap year;
\pcalendardatetojulian{1403-12-29+2}\jd\pcalendarjuliantodate\jd\y\m\d%
plus two days is \iso{\number\y-\number\m-\number\d}.

\paragraph{Counts used internally.} \cs{pcalendardatetojulian} uses \cs{count1}, \cs{count2},
\cs{count3} and \cs{count8} as scratch registers. Do not keep values in those.

% ---------------------------------------------------------------------------
\section{Names}\label{sec:names6}

\macro|\pcalendarmonthname{n}|
Persian month name, \texttt{n} from 1 to 12. An argument outside this range prints nothing.
\ex|\pcalendarmonthname{1}|{\pcalendarmonthname{1}}
\gap
All twelve, in order:\gap
\gap
\begin{tabular}{llll}
\MROW{1}\\
\MROW{2}\\
\MROW{3}
\end{tabular}

\macro|\pcalendargregorianmonthname{n}|
Persian transliteration of the Gregorian month name, \texttt{n} from 1 to 12.
\ex|\pcalendargregorianmonthname{3}|{\pcalendargregorianmonthname{3}}

\macro|\pcalendarweekdayname{n}|
Weekday name, \texttt{n} from 0 to 6, where \textbf{0 is Saturday} (the first day of the Persian
week) and 6 is Friday.
\gap
\foreach \w in {0,...,6}{\texttt{\w}~\pcalendarweekdayname{\w}\quad}

\macro|\pcalendarweekdayshortname{n}|
One-letter weekday abbreviation, same numbering.
\gap
\foreach \w in {0,...,6}{\texttt{\w}~\pcalendarweekdayshortname{\w}\quad}

\macro|\pcalendaryearname{year}|
Name of a Jalali year in the twelve-animal cycle. \texttt{year} must be between 0 and 2099;
other values print nothing.
\ex|\pcalendaryearname{1403}|{\pcalendaryearname{1403}}
\gap
The twelve animals in order, starting from year 1380:\gap
\gap
\begin{tabular}{llll}
\YROW{1380}\\
\YROW{1381}\\
\YROW{1382}
\end{tabular}

\paragraph{The weekday of a date.} The package has no weekday macro, but the day number gives you
one. \pkg{pgfcalendar} numbers Monday as 0, and \cs{pcalendarweekdayname} numbers Saturday as 0, so
shift by two:
\begin{verbatim}
\newcount\jd \newcount\wd
\pcalendardatetojulian{1403-01-01}\jd
\pgfcalendarjuliantoweekday{\jd}{\wd}
\advance\wd by 2 \ifnum\wd>6 \advance\wd by -7 \fi
\pcalendarweekdayname{\wd}
\end{verbatim}
\pcalendardatetojulian{1403-01-01}\jd%
\pgfcalendarjuliantoweekday{\jd}{\wd}%
\advance\wd by 2 \ifnum\wd>6 \advance\wd by -7 \fi%
gives \pcalendarweekdayname{\wd} (1 Farvardin 1403 was a Wednesday).

% ---------------------------------------------------------------------------
\section{Recipes}\label{sec:recipes7}

\paragraph{Days between two Jalali dates.}
\begin{verbatim}
\newcount\a \newcount\b
\pcalendardatetojulian{1403-01-01}\a
\pcalendardatetojulian{1404-01-01}\b
\advance\b by -\a
\number\b
\end{verbatim}
\pcalendardatetojulian{1403-01-01}\jd\newcount\ja\ja=\jd
\pcalendardatetojulian{1404-01-01}\jd\advance\jd by -\ja
prints \number\jd{} days (1403 is a leap year).

\paragraph{A list of consecutive days.} Typeset the 5 days after 1 Farvardin 1403 with a loop. Each date
is converted and printed right away, because the result macros are overwritten by the next
conversion:
\begin{verbatim}
\newcount\jd
\foreach \x in {0,...,4}{%
  \pcalendardatetojulian{1403-01-01+\x}\jd
  \pcalendarjuliantodate\jd\y\m\d
  \number\y-\number\m-\number\d\par
}
\end{verbatim}
\foreach \x in {0,...,4}{%
  \pcalendardatetojulian{1403-01-01+\x}\jd
  \pcalendarjuliantodate\jd\y\m\d
  \par\noindent\hspace*{1em}\iso{\number\y-\number\m-\number\d}%
}
\par\medskip
Using \texttt{\cs{x}} inside \cs{pcalendardatetojulian} works, because that macro expands its argument.

\paragraph{Declare registers once.} \cs{newcount} inside a loop allocates a new register at every
iteration and eventually fails with ``No room for a new \cs{count}''. Declare your counters before the
loop.

% ---------------------------------------------------------------------------
\section{How it works}\label{sec:how8}

The conversions are plain TeX register arithmetic, with no external program and no Lua. Two internal
macros do all the work. Every public macro calls one of them and then prints or copies the result
registers.

\begin{description}
  \item[\texttt{\textbackslash pcal@gregoriantojalali\{Y-M-D\}}] Gregorian to Jalali. It first decides
    whether the Gregorian year is a leap year, builds the month-length tables of both calendars, and
    counts the days from the start of the Gregorian year to the date. It then walks the Jalali month
    lengths until that day count is used up. The Jalali leap year rule is the common
    \emph{33-year cycle}: eight leap years in every 33.
  \item[\texttt{\textbackslash pcal@jalalitogregorian\{Y-M-D\}}] Jalali to Gregorian. It computes the
    number of days since a fixed epoch from the year (365 days per year, plus 8 leap days per 33 years,
    plus the leap days of the partial cycle), adds the days of the preceding months, and then splits
    that total into 400-year, 100-year, 4-year and 1-year Gregorian blocks.
\end{description}

\paragraph{Result registers.} Jalali results are in \cs{pcal@y}, \cs{pcal@m}, \cs{pcal@d}. Gregorian
results are in \cs{pcal@gy}, \cs{pcal@gm}, \cs{pcal@gd}. The leap flags are \cs{ifpcal@leap}
(Gregorian) and \cs{ifpcal@kabiseh} (Jalali). The month lengths are held in registers whose names end in
roman numerals (\cs{pcal@latini} to \cs{pcal@latinxii} and \cs{pcal@persiani} to
\cs{pcal@persianxiii}). Dey appears twice: the first slot is the part of Dey left in the Gregorian year, the last is
the whole month, which is why there are thirteen.

\paragraph{Julian day bridge.} \cs{pcalendardatetojulian} converts the Jalali date to Gregorian with the
Jalali-to-Gregorian macro above and hands the result to \cs{pgfcalendardatetojulian}.
\cs{pcalendarjuliantodate} does the reverse: \cs{pgfcalendarjuliantodate} gives a Gregorian date, which the
Gregorian-to-Jalali macro turns into the Jalali date.

\paragraph{Year names.} The twelve animals repeat every 12 years. A second, three-step shift is applied
once per century, so that year 0 starts the cycle at a different animal than year 100. This is why the
macro stops at 2099.

\paragraph{No register leaks.} All registers are allocated once when the package loads, so any number of
conversions can be made in one document. A test in the repository (\texttt{leaks.lvt}) checks this.

% ---------------------------------------------------------------------------
\section{Limitations and known issues}\label{sec:limitations9}

\begin{itemize}
  \item \textbf{No input checking.} A month 13, day 40 or malformed string gives wrong output or a TeX
    error, not a package message.
  \item \textbf{Literal dates only.} The date argument of \cs{pdate}, \cs{gdate}, \cs{pisodate} and
    \cs{gisodate} is split at the hyphens as it stands; a macro holding a date is not expanded. Write
    \cs{pdate}\texttt{\{2015-03-21\}}, not \cs{pdate}\cs{mydate}. (\cs{pcalendardatetojulian} does
    expand its date.)
  \item \textbf{Approximate Jalali leap rule.} The 33-year rule is an arithmetic approximation of the
    official calendar, which is defined by astronomical observation. The two agree for the years in
    ordinary use, but not necessarily for all years. For distant years, check against an authoritative
    table.
  \item \textbf{A known crash.} \cs{pcal@jalalitogregorian}\texttt{\{1279-10-11\}} stops with ``Missing
    number, treated as zero''; neighbouring dates work. The cause has not been found.
  \item \textbf{Nested loops.} The conversions use \cs{loop} from the \LaTeX{} kernel internally. A
    \cs{loop}\ldots\cs{repeat} of your own that calls them must run its body inside a group
    (\cs{begingroup}\ldots\cs{endgroup}), otherwise it stops after one pass. \cs{foreach} is not affected.
  \item \textbf{Shared registers.} Every conversion overwrites the result registers, as described in
    section~\ref{sec:day5}.
  \item \textbf{Year names} work for 0 to 2099 only.
\end{itemize}

% ---------------------------------------------------------------------------
\section{Version history}\label{sec:version11}

\begin{description}
  \item[0.51 (2026-11-01)] Documentation only: the manuals now say that only printing the Persian names
    needs a Unicode engine and a Persian font; the Persian manual was rewritten.
  \item[0.5 (2026-11-01)] Fixed the register leak (several \cs{newcount} per call) and the undeclared
    month-length register; added this manual and
    \cs{pcalendardate}/\cs{pcalendarversion}.
  \item[0.4 (2023-11-07)] Ready for CTAN.
  \item[0.3 (2015-08-14)] Animal name of a Jalali year, 0 to 2099.
  \item[0.2 (2015-08-09)] Jalali to Gregorian conversion.
  \item[0.1 (2015-07-25)] Initial version.
\end{description}

% ---------------------------------------------------------------------------
\section{License}\label{sec:license12}

This work may be distributed and/or modified under
the conditions of the \LaTeX{} Project Public License, either version 1.3c of this license or (at your
option) any later version. The latest version is at \texttt{http://www.latex-project.org/lppl.txt}. This
work has the LPPL maintenance status ``author-maintained''.

\end{document}
